\section{The verdict algebra: Mizrahi's impossibility as a type-level invariant}
\label{sec:verdict}

This is the object the paper is built around. Everything before it is method
(\Cited{}) and stimulus (\Modelled{}); here is the one \Proven{} contribution---a
scoring algebra whose \emph{shape} makes a certain dishonest outcome impossible to
express. The design goal is blunt: it must be structurally impossible for the
benchmark to credit any monitor with detecting a fault that is provably
undetectable. We achieve this not by a runtime check that could be bypassed, but by
the enumeration of possible verdicts itself.

\subsection{The four-value \code{Verdict} and the absent fifth}
\label{sec:verdictenum}

The scorecard runs a monitor---a function $\PL(k)$ giving the protection level it
declares at each epoch from the epoch's observables, with the true error held
hidden---against every scenario in a menu and emits, per scenario, one of exactly
four verdicts:
\begin{equation}
  \code{Verdict} \;\in\; \big\{\,
  \CovPass{},\;\;
  \CovFail{},\;\;
  \Absorbed{},\;\;
  \Unabsorbed{}
  \,\big\}.
  \label{eq:verdict}
\end{equation}
There is no fifth value. In particular there is no \textsf{Detected} and no
\textsf{DetectionMissed}: the enumeration simply does not contain a token for
``the monitor detected the fault.'' This absence is the invariant.

The two \code{Coverage} verdicts apply only to \code{Detectable} scenarios and are
the coverage decision of Section~\ref{sec:coverage}: \CovPass{} when the empirical
\HMI{} rate is within the monitor's stated integrity risk, \CovFail{} otherwise.
The two \code{Undetectable} verdicts apply only to \code{Undetectable} scenarios
and are an \emph{absorption} decision:
\begin{equation}
  \text{undetectable scenario} \;\Longrightarrow\;
  \begin{cases}
    \Absorbed{}, & \PL(k)\ \ge\ |e(k)| \ \ \text{for every epoch } k,\\[2pt]
    \Unabsorbed{}, & \text{otherwise.}
  \end{cases}
  \label{eq:absorb}
\end{equation}
The dispatch is on the scenario's cited detectability flag
(Section~\ref{sec:detectability}), so the mapping from fault kind to verdict kind
is fixed before any monitor is run.

\subsection{The impossibility theorem}
\label{sec:impossibility}

\begin{theorem}[Undetectable faults are never reported as detected]
\label{thm:invariant}
For every monitor $\PL(\cdot)$, every alert limit, every stated integrity risk,
and every scenario flagged \code{Undetectable}, the emitted verdict is either
\Absorbed{} or \Unabsorbed{}. No execution path of the scorecard emits a
detection-flavoured verdict for such a scenario, and no such verdict exists to be
emitted.
\end{theorem}

\begin{proof}
The verdict set \eqref{eq:verdict} contains no detection token, so no value
denoting detection can be constructed. The scorecard's dispatch branches on the
scenario's \code{Detectability}: the \code{Undetectable} branch computes only the
absorption predicate \eqref{eq:absorb} and can return only \Absorbed{} or
\Unabsorbed{}; the \code{Coverage} verdicts are unreachable from that branch. Since
the two branches partition the scenarios and the \code{Undetectable} branch's range
is exactly $\{\Absorbed{},\Unabsorbed{}\}$, the conclusion follows by exhaustion of
the (finite, four-valued) codomain.
\end{proof}

The proof is deliberately short, and its shortness is the point: the theorem holds
by the \emph{shape of the type}, not by a defensive conditional a future refactor
could remove. Mizrahi's impossibility \citep{mizrahi7384}---that a symmetric
single-path delay cannot be detected by any protocol---is thus encoded not as prose
in a comment but as an unrepresentable state. The crate's tests make the invariant
executable: \code{undetectable\_never\_yields\_a\_coverage\_verdict} runs an
undetectable scenario across protection levels spanning zero to $10^{6}$ and
asserts the verdict is always an absorption verdict, and the honesty-gate test
\code{undetectable\_faults\_never\_described\_as\_detected} scans the benchmark
source so that no line naming an undetectable class may also assert detection
(negations such as ``never detected'' excepted). The type makes the dishonest
outcome unconstructable; the tests make sure the prose around it stays honest too.

\subsection{Absorption as a positive scoring rule}
\label{sec:absorptionrule}

The elegance of \eqref{eq:absorb} is that it converts a negative result into a
positive incentive. Because an undetectable offset cannot be detected, the only
honest question a benchmark can ask about it is: \emph{was the monitor's protection
level large enough to carry it?} If $\PL(k)\ge|e(k)|$ at every epoch, the
undetectable offset never escaped the declared bound---the monitor was
\emph{honest about its own uncertainty}, having reserved enough protection level to
cover a fault it could never have seen. That is \Absorbed{}, and it is the correct
passing outcome. If the protection level ever fell below the offset, the monitor
under-reserved against the undetectable, and that is \Unabsorbed{}. Note what is
\emph{not} rewarded: nowhere does a monitor gain credit for flagging, alerting on,
or detecting the symmetric delay, because there is no verdict for that and no
adversary it could distinguish. A monitor is rewarded strictly for carrying
sufficient protection level---exactly the behaviour Mizrahi's result says is the
only available defence.

\subsection{The aggregate scorecard}
\label{sec:aggregate}

Across a menu, the scorecard aggregates two honesty-relevant totals: the total
\HMI{} count $n_{\HMI{}}^{\text{total}}$ summed over all scenarios (the count of
hazardous undetected-error epochs, the benchmark's headline integrity number), and
a boolean $\mathrm{all\_undetectable\_absorbed}$ that is true iff every
undetectable scenario earned \Absorbed{}. The first measures how a monitor fares on
the detectable menu; the second measures whether it honestly reserved against the
undetectable. Neither is an accuracy claim about the monitor---the benchmark still
asserts nothing about how good the monitor is in absolute terms---and both are pure
tallies of the fixed scoring rule. A monitor with a small
$n_{\HMI{}}^{\text{total}}$ and $\mathrm{all\_undetectable\_absorbed}=\text{true}$
has covered the detectable faults and honestly carried the undetectable ones; a
monitor that fails either is exposed, on the same menu, by the same rule. The
demonstration of Section~\ref{sec:demo} reads exactly these two totals off a real
reference monitor and a broken one.
